\chapter{Flight phases}
\label{impl:phases}
The implementation of Chapter~\ref{ch:phases} is deliberately separate from preprocessing:
it reads the final \path{fixes.parquet} table and writes all inferred quantities below
\path{derived/segmentation/}. Consequently a phase-analysis rerun cannot alter the
trajectory record on which the global transport chapter was measured.

\section{Segmentation into flight phases}
\label{impl:segmentation}
The package \path{soaring.analysis.segmentation} has five responsibilities. Its feature
builder groups the streamed fix table by $(\texttt{source},\texttt{flight\_id},
\texttt{segment\_id})$, refuses a segment with native cadence above \SI{10}{\second},
and evaluates the continuous features of Sec.~\ref{sec:phasefeatures} in physical-time
windows. It carries the preprocessing \texttt{z\_reconstructed} and \texttt{edge} flags:
a window touching either is non-classifiable and starts a new valid block. The HMM
therefore receives one sequence per retained valid block; its \texttt{lengths} argument
prevents a fictitious state transition at every preprocessing, quality or sampling join.

\path{configs/segmentation.yaml} is the single record of the decision grid, feature
window, state number, full covariance emission model, split fractions, seeds, restart
count and fitting-memory caps. The model artefact contains the fitted \texttt{GaussianHMM},
the training-only standardizer, the numerical-state-to-phase mapping and a JSON copy of
that configuration. A deterministic flight-level manifest makes the train/validation/test
partition reproducible independently of the order in which Parquet row groups were stored;
a second fitting-sample manifest records the exact whole valid blocks, their order and
their observation counts. Blocks beyond the configured memory limits are omitted rather
than cut at an artificial transition.

Manual labels are CSV or Parquet intervals with columns
\texttt{source}, \texttt{flight\_id}, \texttt{segment\_id}, \texttt{t\_start},
\texttt{t\_end}, \texttt{state}, \texttt{split}, and \texttt{annotator}; the time
bounds use the half-open convention $[\texttt{t\_start},\texttt{t\_end})$. The validator
refuses missing fields, invalid state names, invalid split names and overlapping intervals.
Train labels resolve HMM label switching with a one-to-one Hungarian assignment; the fit
itself remains unsupervised.

The command-line entry point is, for example,
\begin{verbatim}
uv run python scripts/segment_flights.py train --discipline paragliders
uv run python scripts/segment_flights.py apply --discipline paragliders
\end{verbatim}
The two-model rule is enforced by running the command once per discipline: model, scaler
and transition matrix are never shared by paragliders and hang gliders.  With no labels,
\texttt{train} records a provisional emission-signature mapping so that the archive can
be decoded and a blinded annotation set prepared.  Once train intervals exist,
\texttt{calibrate} replaces that permutation with the Hungarian assignment and atomically
remaps the existing point and run tables; it does not repeat the expensive preprocessing
scan or Baum--Welch fit.  \texttt{evaluate} then reads validation or test labels without
modifying the model.

The ten HMM restarts are independent processes and the full archive decoder likewise
distributes flights over at most eight worker processes, as fixed in the YAML.  Each fit
worker is restricted to one BLAS thread, preventing ten processes from each starting a
second pool of threads.  Work is exchanged as bounded batches of whole flights, so this
parallelization changes neither their order in the output nor any sequence boundary.

\section{Phase-resolved observables}
\label{impl:phaseobs}
Decoding writes \path{phase_points.parquet}, one row per \SI{10}{\second} decision time,
with the four feature values, the Viterbi state and its three smoothed posterior
probabilities. A feature-window edge or \texttt{quality\_masked} interval is written
explicitly as \texttt{unclassified} rather than padded or guessed.
\path{phase_segments.parquet} reduces consecutive equal Viterbi
states to runs, retaining their time span, duration, number of decision points, mean
posterior confidence and left/right censoring flags. Those flags must accompany all
duration-law fits, since a phase that reaches an observable edge is not known to have ended.
The companion \path{phase_coverage.parquet} has one row per preprocessing segment and
records whether it was decoded, lacked any safe centred window, or was skipped because its
native cadence exceeded \SI{10}{\second}; consequently an excluded segment cannot disappear
silently from a coverage calculation.

\section{Reproducing the reference results}
\label{impl:reproducing}
The evaluation routine aligns only manually labelled intervals with decoded decision points.
It reports accuracy, state-wise precision/recall/F1 and macro-F1, with a 95\% bootstrap
interval obtained by resampling whole flights rather than correlated grid points. Validation
and test reports are separate JSON artefacts. The test set is not read by fitting, state
mapping or model selection. The Chapter-4 reporter additionally writes likelihood per
observation, occupancy and uncensored mean run duration for train, validation and test, all restart
scores and convergence flags, transition and emission-covariance diagnostics, example
decoded trajectories, confusion matrices and the generated metric table.

\section{Modelling and simulation}
\label{impl:modelling}
The phase outputs are the input contract for the waiting-time, glide-step and phase-memory
measurements of Sec.~\ref{sec:phaseobs}; the latter remain separate analyses. In particular,
the transition matrix and the run table make the geometric HMM dwell-time prediction
testable rather than assumed.

\section{Tooling}
\label{impl:tooling}
The package tests cover physical-time feature construction, cadence invariance, refusal to
cross a segment boundary, one-to-one state mapping, split stability, run censoring, model
artefact persistence and HMM decoding. The operator-facing annotation and output contract
is also documented in \path{docs/guide/flight-phase-segmentation.md}.

The trajectory viewer exposes the same decoder through a \emph{Flight phase (HMM)} colour
mode.  It intentionally does not scan the archive-scale \path{phase_points.parquet} on
each click: it caches the small model artefact and applies the common-grid feature builder
and Viterbi decoder only to the selected flight already in memory.  A separate display-run
identifier changes at every state, time-gap and preprocessing-segment boundary, preventing
the renderer from joining two non-contiguous occurrences of the same coloured phase.  The
plot title and status line distinguish provisional emission-signature names from a manual train-set
calibration.

\stopcontents[impl]
